\uuid{ZHqe}
\titre{Tangentes}
\auteur{hanani ; mzari}
\organisation{exo7}
\licence{CC BY-NC-SA 4.0}
\datecreate{2020-01-20}
\niveau{L1}
\module{Géométrie}
\chapitre{Courbes planes}
\sousChapitre{Courbes paramétrées}
\difficulte{2}
\theme{Courbes paramétrées, tangentes}
\exo7qcmid{lille-2.courbes.003}

\contenu{
\begin{question}
\qclassification{244.01}
\qtags{motcle=tangentes}
Un avion en papier a effectué un vol suivant la trajectoire $\Gamma$ donnée par
$$x=t-2\sin t\quad \mbox{et} \quad y=4-3\cos t.$$
\begin{answers}
\bad{A l'instant $t=\pi$, l'avion volait en position verticale.}
\good{A l'instant $t=\pi$, l'avion volait en position horizontale.}
\bad{A l'instant $t=\pi/2$, l'avion volait suivant la droite d'équation $y=3x+10$.}
\good{A l'instant $t=\pi/3$, l'avion volait en position verticale.}
\end{answers}
\begin{explanations}
On a : $y'(\pi)=0$ et $x'(\pi)\neq 0$. Donc la tangente au point de paramètre $\pi$ est horizontale. A l'instant $t=\pi/2$, l'avion volait suivant la tangente \`a $\Gamma$ à cet instant. C'est-à-dire suivant la droite d'équation $y=3x+10-3\pi /2$. Enfin, $x'(\pi/3)=0$ et $y'(\pi/3)\neq 0$. Donc la tangente au point de paramètre $\pi/3$ est verticale.
\end{explanations}
\end{question}
}
